% Guide-only presentation. Packages and shared styles come from Packages/.
% This file is optional: its contents could be placed directly in dev.tex.
\geometry{top=27mm,bottom=26mm,left=27mm,right=27mm,
  headheight=24pt,footskip=13mm}
\hypersetup{hidelinks,pdftitle={Package Reference},
  pdfsubject={LaTeX package and command reference},bookmarksopen=true}
\setlength{\parindent}{0pt}
\setlength{\parskip}{4pt}
\setlength{\emergencystretch}{2em}
\renewcommand{\arraystretch}{1.15}
\setcounter{tocdepth}{2}
\raggedbottom
\fancyhf{}
\fancyhead[R]{\includegraphics[height=7mm]{Images/logo_header.png}}
\fancyfoot[C]{\thepage}
\renewcommand{\headrulewidth}{0.4pt}
\renewcommand{\footrulewidth}{0.4pt}
\pagestyle{fancy}
\fancypagestyle{plain}[fancy]{}

% Small helpers for consistent headings and examples; no new listing styles.
\newcommand{\packagegroup}[2]{\clearpage\section{#2}\label{class:#1}}
\newcommand{\packageentry}[2]{\Needspace{9\baselineskip}\subsection{#2}\label{#1}}
\newcommand{\partlabel}[1]{\par\smallskip\noindent\textbf{#1}\par\nobreak}
% Distinct environment name lets examples contain literal lstlisting blocks.
\lstnewenvironment{referencecode}{\lstset{style=latex,
  breaklines=false,upquote=true,rulecolor=\color{black!30}}}{}
